\subsection{Benchmarks Existentes}

La evaluación sistemática de las capacidades de los LLMs para utilizar
herramientas y completar tareas complejas ha motivado el desarrollo de
múltiples benchmarks especializados. Esta sección revisa los más relevantes,
organizados en dos categorías: benchmarks de tool-use (que evalúan la capacidad
de invocar herramientas correctamente) y benchmarks de agentes (que evalúan la
capacidad de completar tareas end-to-end). La sección concluye identificando el
vacío que esta investigación busca llenar.

\subsubsection{Benchmarks de tool-use}

Los benchmarks de tool-use evalúan la capacidad de un LLM para seleccionar
herramientas apropiadas, generar parámetros correctos y encadenar múltiples
llamadas cuando es necesario.

\paragraph{ToolBench.} \citet{toolbench2023} presentan un benchmark a gran
escala con más de 16,000 APIs del mundo real organizadas en categorías
funcionales. ToolBench evalúa no solo la selección de herramientas individuales,
sino la capacidad del modelo para resolver tareas que requieren múltiples
llamadas a APIs heterogéneas. Los autores proponen un método de planificación
basado en árboles de decisión (DFSDT) que mejora significativamente el
rendimiento en cadenas de herramientas complejas. La principal contribución de
ToolBench es demostrar que la escala del espacio de herramientas ---cientos o
miles de APIs disponibles--- es un factor crítico en el rendimiento de los
modelos.

\paragraph{API-Bank.} \citet{apibank2023} construyen un benchmark para evaluar
LLMs aumentados con herramientas a través de tres niveles de dificultad: Nivel
1 evalúa la invocación correcta de una API individual; Nivel 2 requiere
recuperar la API apropiada de un conjunto amplio; Nivel 3 demanda planificación
y ejecución de múltiples APIs con dependencias entre ellas. API-Bank
proporciona un entorno de ejecución donde las llamadas a APIs se ejecutan
realmente, permitiendo evaluar no solo la generación de parámetros sino la
corrección funcional de las respuestas.

\paragraph{Berkeley Function Calling Leaderboard (BFCL).} El BFCL
\citep{bfcl2024} estandariza la evaluación de function calling a través de
múltiples categorías: llamadas simples, llamadas múltiples, llamadas paralelas,
selección de funciones relevantes de un conjunto amplio, y detección de
consultas que no requieren ninguna herramienta. Su contribución principal es la
\textbf{estandarización}: al proveer un leaderboard público con métricas
uniformes, permite comparar directamente la capacidad de function calling entre
modelos de diferentes proveedores bajo condiciones idénticas.

\subsubsection{Benchmarks de agentes}

A diferencia de los benchmarks de tool-use, que evalúan interacciones aisladas
con herramientas, los benchmarks de agentes evalúan la capacidad de completar
tareas complejas que requieren múltiples pasos de razonamiento, planificación y
ejecución.

\paragraph{SWE-Bench.} \citet{swebench2024} presentan un benchmark que evalúa
la capacidad de agentes para resolver issues reales de repositorios de código
abierto en GitHub. Cada instancia consiste en una descripción de un bug o
feature request y el estado del repositorio correspondiente; el agente debe
producir un patch que resuelva el issue y pase los tests del repositorio. SWE-Bench
es particularmente relevante porque evalúa una cadena completa de capacidades:
comprensión de código existente, razonamiento sobre el problema, generación de
código y validación mediante tests. Los resultados reportados muestran que
incluso los mejores agentes resuelven menos del 30\% de las instancias,
evidenciando la dificultad de las tareas de ingeniería de software autónoma.

\paragraph{WebArena.} \citet{webarena2024} construyen un entorno realista de
navegación web con sitios funcionales (e-commerce, foros, herramientas de
gestión de contenido) donde los agentes deben completar tareas complejas
mediante interacciones con la interfaz web. WebArena evalúa la capacidad del
agente para planificar secuencias de acciones en un entorno dinámico donde el
estado cambia con cada interacción, un escenario análogo al uso de
herramientas MCP donde cada llamada puede alterar el estado del sistema.

\paragraph{GAIA.} \citet{gaia2024} proponen un benchmark para evaluar
asistentes de IA generales en tareas que requieren razonamiento multi-paso,
uso de herramientas web y manejo de archivos multimodales. Las tareas de GAIA
están diseñadas para ser fáciles de verificar (las respuestas son factuales y
unívocas) pero difíciles de resolver, requiriendo que el agente combine
múltiples capacidades. GAIA es relevante para esta investigación porque sus
tareas implican frecuentemente cadenas de herramientas ---exactamente el
escenario donde la ejecución de código puede ofrecer ventajas sobre el tool
calling convencional.

\subsubsection{El vacío identificado}

Los benchmarks revisados comparten una característica en común: todos evalúan
\textit{qué} logra el agente (correctitud de la respuesta, resolución del
issue, completitud de la tarea), pero ninguno evalúa \textit{cómo} el agente
ejecuta las herramientas. Específicamente, se identifican tres vacíos en la
literatura:

\begin{enumerate}
    \item \textbf{Modo de ejecución como variable}: ningún benchmark existente
    compara el rendimiento de un mismo agente utilizando tool calling
    convencional versus ejecución de código como modos alternativos de
    interacción con herramientas.

    \item \textbf{Eficiencia de tokens}: los benchmarks existentes reportan
    métricas de éxito/fracaso, pero no miden sistemáticamente el consumo de
    tokens como función del modo de ejecución. Dado que el costo operativo de
    los agentes es directamente proporcional al consumo de tokens, esta métrica
    es crítica para la viabilidad de despliegues en producción.

    \item \textbf{Complejidad de herramientas como variable}: no existe un
    benchmark que varíe sistemáticamente la complejidad del servidor de
    herramientas (número de herramientas, grado de concatenación requerido) para
    evaluar cómo escala cada modo de ejecución.
\end{enumerate}

Esta tesis aborda estos vacíos mediante un diseño experimental factorial $2
\times 2$ (modo de ejecución $\times$ número de agentes) evaluado sobre
servidores MCP de complejidad creciente, midiendo simultáneamente correctitud,
consumo de tokens y latencia. El diseño experimental se detalla en el Capítulo
3.
